\section{Affine common ideals with a shared tail give an original-support Green form}
\label{sec:r11-affine-shared-tail}

\paragraph{Status and construction.}
\textbf{PROVED THIS ROUND}, subject to exact-source independent review.
The construction contributor is \texttt{/root/r11\_challenge}.
The new input is the explicit two-chart analytic block theorem GB,
not an asserted factorization of an arbitrary analytic matrix.
The conormal calculation used to find the invertibility neighborhood
is the elementary frozen ICC calculation, recalled below.
Neither the cancellation locus nor its closure is assumed separated
from the remaining support in $X$.

\subsection{An algebraic certificate on one affine neighborhood}

Keep the original finite $K/\mathbf Q_p$, its absolute value, and a
smooth projective geometrically integral $X/K$ of dimension $n$.
Let $n\ge r\ge2$, let $L$ be one line bundle with one of the actual
metrics allowed in frozen CI, and let
\[
 s_0=(f_1,f_2,f_3,\ldots,f_r),\qquad
 s_1=(g_1,g_2,f_3,\ldots,f_r)
 \tag{AST1}
\]
be regular global $r$-tuples of sections of $L$.  The last $r-2$
sections are literally identical, not merely generators of equal
tail ideals.  For $r=2$ the tail is empty.  Put
\[
 \begin{split}
 Y_i&=V(s_i),\\
 Z&=[Y_1]-[Y_0],\\
 U&=X^{\rm an}\setminus|Z|,\\
 W_0&=X^{\rm an}\setminus(|Y_0|\cup|Y_1|),\\
 C_U&=(|Y_0|\cup|Y_1|)\cap U.
 \end{split}
 \tag{AST2}
\]
All cycles, multiplicities and supports in (AST2) are the original
ones.  In particular $C_U$ can accumulate at $|Z|$ in $X^{\rm an}$.

Supply a \emph{single Zariski affine open} $V_{\rm alg}\subset X$
such that
\[
 C_U\subset V_{\rm alg}^{\rm an},\qquad
 \mathcal I_{Y_0}|_{V_{\rm alg}}
       =\mathcal I_{Y_1}|_{V_{\rm alg}}=:\mathcal I.
 \tag{AST3}
\]
This is an independently checkable additional hypothesis.  Equality
of cycle multiplicities alone is not (AST3).  No trivialization of
$L|_{V_{\rm alg}}$ is required.

\begin{lemma}[A global block matrix over the original affine ring]
There are matrices $A\in M_2(\mathcal O(V_{\rm alg}))$ and
$B\in M_{2,r-2}(\mathcal O(V_{\rm alg}))$ such that
\[
 s_1=T s_0,\qquad
 T=\begin{pmatrix}A&B\\0&I_{r-2}\end{pmatrix}
 \tag{AST4}
\]
on all of $V_{\rm alg}$.  The regular function $\det A$ has no
zero on the underlying common zero scheme in $V_{\rm alg}$.
The ordinary open
\[
 N=\bigl(V_{\rm alg}\setminus V(\det A)\bigr)^{\rm an}\cap U
 \tag{AST5}
\]
therefore contains $C_U$, and $T$ is invertible on all of $N$.
The common zero scheme need not be reduced.
\end{lemma}

\begin{proof}
The first tuple defines a surjection of quasi-coherent sheaves
\[
             \mathcal O_{V_{\rm alg}}^r
                       \longrightarrow \mathcal I\otimes L.
 \tag{AST6}
\]
Because $V_{\rm alg}$ is affine, its map of global sections is
surjective.  Explicitly, under the equivalence with modules over
$\mathcal O(V_{\rm alg})$, the cokernel of the module map has
zero localization at every prime, hence is zero.  This applies
to the quasi-coherent kernel and cokernel of (AST6); it does not
require the invertible module corresponding to $L$ to be free.
The sections $g_1,g_2$ belong to
$\Gamma(V_{\rm alg},\mathcal I\otimes L)$ by (AST3).
Lift them by (AST6), giving the first two rows of $T$.
The last rows in (AST4) are the displayed identity because
the tail in (AST1) is exactly shared.

At a local point of the common zero scheme, trivialize $L$ and
write the tuples as regular-sequence generators $f,g$ of the
same ideal $I$ in the local ring.  Locally there is a matrix
$S$ with $f=Sg$.  Thus $(ST-I_r)f=0$.  The frozen ICC
regular-sequence relation calculation says that each coefficient
of every relation among the $f_i$ is in $I$.  Consequently
\[
                     ST\equiv I_r\pmod I.
 \tag{AST7}
\]
Taking determinants over the commutative quotient ring gives
that $\det T=\det A$ is a unit modulo $I$.  In particular
it is nonzero at every point of the underlying zero scheme.
No passage to its reduction was used.  A common scalar change
of frame of $L$ leaves the coefficient matrix (AST4) unchanged.

Since (AST3) gives the same two zero schemes on $V_{\rm alg}$,
every point of $C_U$ belongs to that common scheme and to the
nonvanishing locus of $\det A$.  Formula (AST5) is an ordinary
analytic open, lies in the original $U$, and contains $C_U$.
The adjugate formula makes $A^{-1}$ and $T^{-1}$ analytic there.
All coefficients and inverses are over the original $K$.
\end{proof}

\subsection{Actual relative primitives and the all-test identity}

On the entire open $N$, the matrix in (AST4) has the concrete
finite decomposition
\[
 T=\begin{pmatrix}I_2&B\\0&I_{r-2}\end{pmatrix}
                    \begin{pmatrix}A&0\\0&I_{r-2}\end{pmatrix},
 \qquad
 \begin{pmatrix}I_2&B\\0&I_{r-2}\end{pmatrix}
       =\prod_{i=1}^{2}\prod_{j=3}^{r} E_{ij}(B_{i,j-2}).
 \tag{AST8}
\]
Fix, for example, lexicographic order in the finite product.
The factors commute because their targets are in $\{1,2\}$
and their sources are in $\{3,\ldots,r\}$; no product of two
of their off-diagonal matrix units is nonzero.  The product is
empty when $r=2$.  Formula (AST8) is a direct multiplication
identity for this block shape, not a general matrix factorization
result.  The analytic entries of $B$ may vanish or be unbounded.

Apply GB23 on $N$ to the first executed block $A\oplus I$.
It constructs a whole-neighborhood relative actual pair by
two ordinary Gaussian charts and the explicitly proved native
overlap relation.  Then execute each shear in (AST8), using
the full native SHT pair, including its coefficient-zero
extension.  Evaluate each pair at the actual intermediate
tuple, and add the resulting finitely many pairs and residues.
All intermediate matrices are invertible on $N$, so all
intermediate tuples have the same common zero set there.
This gives actual forms
\[
 \begin{split}
 p=(p',p'')&\quad\hbox{on }N\cap W_0,\\
 E&\in B^{r-1,r-1}(N),
 \end{split}
 \qquad
 G_1-G_0=d'p'+d''p''+E|_{N\cap W_0}.
 \tag{AST9}
\]
Here the only nonzero extending residue is the one explicitly
constructed in GB22 for the block; each shear has determinant
one and contributes no curvature residue.  This assertion
does not discard the differentiated partition terms in GB22.

\begin{theorem}[Affine shared-tail cancellation]
Under (AST1)--(AST3), the original cycle $Z$ admits an actual
\[
                       g_Z\in B^{r-1,r-1}(U)
 \tag{AST10}
\]
over the original field and absolute value such that
\[
 \int_U g_Z\wedge d'd''\eta=-\int_Z\eta
       \quad\hbox{for every }
       \eta\in B_c^{n-r,n-r}(X^{\rm an}).
 \tag{AST11}
\]
The theorem holds in every $n\ge r\ge2$.  It imposes no
closedness of $C_U$ in $X^{\rm an}$ and no separation of its
closure from $|Z|$.
\end{theorem}

\begin{proof}
Use (AST9) and the frozen RSP cutoff on $U$: choose an
ordinary smooth actual $\chi$ with support closed in $U$
and contained in $N$, equal to one on an ordinary neighborhood
of $C_U$.  With the RSP closed-support zero extensions, put
\[
       g_Z=G_1-G_0-d'(\chi p')-d''(\chi p'')
                                \quad\hbox{on }W_0.
 \tag{AST12}
\]
Near $C_U$, (AST9) says that (AST12) equals the actual form
$E$.  The two actual formulas therefore glue on the open
cover of $U$ formed by $W_0$ and that neighborhood of $C_U$.
There is no assertion that the raw difference itself extends.

For any actual test $\eta$ in (AST11), the frozen support
lemma gives $\operatorname{supp}D\eta\Subset W_0$, where
$D=d'd''$.  In particular both correction products with
$D\eta$ in (AST12) have compact support in $W_0$.
Since $d'D\eta=d''D\eta=0$, actual Stokes makes their
integrals zero.  Subtracting the two full actual CI formulas
cancels the identical curvature term from the same metric
on the same $L$, and gives exactly (AST11).  This verifies
all actual compact tests, not only smooth or separated ones.
The use of the cutoff occurs on $U$, never on $X$ across
the remaining support.
\end{proof}

The same frozen support and fixed-presentation seminorm argument
makes $I(g_Z)$ continuous for the full coefficient-final target
topology in degree $(n-r+1,n-r+1)$.  This is the interface of
the actual form just constructed, not a substitute for (AST10).
No claim about a full lower-middle kernel topology or a universal
continuous primitive selection is made.

\subsection{A moving curve in every rank with neither pivot globally available}

The following calibration is over the original $K$ for every
$r\ge2$.  It exhibits nonclosed cancellation and shows that the
two pivot charts in GB can both be needed by this presentation.
It is not a claim that this particular cycle evades every older
certificate.

Take
\[
 X=\mathbf P^1_{[s:t]}\times\mathbf P^r_{[z_0:\cdots:z_r]},
       \qquad L=\mathcal O(1,1),
       \qquad h_j=s z_j+t z_{j-1}\quad(3\le j\le r),
 \tag{AST13}
\]
with any one of the allowed actual metrics on $L$, and set
\[
 \begin{split}
 s_0&=(s z_1,\ s z_2+t z_0,\ h_3,\ldots,h_r),\\
 s_1&=(t z_1-s z_2-t z_0,
          (s-t)z_1+s z_2+t z_0,\ h_3,\ldots,h_r).
 \end{split}
 \tag{AST14}
\]
These are global sections of the same $L$ in every entry.
On $s\ne0$ put $\tau=t/s$.  The exact matrix identity is
\[
 s_1=\left(\begin{pmatrix}\tau&-1\\1-\tau&1\end{pmatrix}
                           \oplus I_{r-2}\right)s_0,
                    \qquad\det A=1.
 \tag{AST15}
\]
The common zero locus over this affine base is the graph
\[
          z_1=0,\qquad z_j=(-\tau)^{j-1}z_0
                     \quad(2\le j\le r).
 \tag{AST16}
\]
Here $z_0$ cannot be zero, because then every projective
coordinate would be zero.  Thus this graph lies in the one
affine open
\[
       V_{\rm alg}=\{s\ne0\}\times\{z_0\ne0\}
                                  \simeq\mathbf A^{r+1}_K.
 \tag{AST17}
\]
Identity (AST15) proves equality of the two ideals throughout
this affine open, with an exactly shared tail.

Let $C$ be the closure of (AST16).  Explicitly it is the graph
of the morphism from $\mathbf P^1$ with projective coordinates
\[
 z_0=s^{r-1},\quad z_1=0,\quad
 z_j=(-1)^{j-1}t^{j-1}s^{r-j}\ (2\le j\le r).
 \tag{AST18}
\]
They have no common zero, and the first projection identifies
$C$ with $\mathbf P^1_K$.  In the fiber $s=0$, set
\[
 \begin{split}
 D_0&=\{s=0,\ z_0=z_2=\cdots=z_{r-1}=0\},\\
 D_1&=\{s=0,\ z_1=z_0,\ z_2=\cdots=z_{r-1}=0\}.
 \end{split}
 \tag{AST19}
\]
An empty list of the equations $z_2,\ldots,z_{r-1}$ is
understood for $r=2$.  Both $D_i$ are projective lines.
They are distinct, and their intersection is the single point
\[
             P_\infty=([0:1],[0:\cdots:0:1]).
 \tag{AST20}
\]
The curve $C$ meets that fiber precisely at $P_\infty$.

\begin{lemma}[Exact schemes and multiplicities in the calibration]
The tuples (AST14) are regular.  Their zero schemes are reduced
local complete intersection curves with cycle decompositions
\[
                     [Y_0]=[C]+[D_0],\qquad
                     [Y_1]=[C]+[D_1].
 \tag{AST21}
\]
There are no further components or embedded associated points.
\end{lemma}

\begin{proof}
On $s\ne0$, successively eliminating $z_1,z_2,\ldots,z_r$
in the chart $z_0\ne0$ proves that $Y_0$ is the smooth graph
(AST16), cut out by a regular sequence.  Formula (AST15)
gives the same scheme for $Y_1$ there.  At $s=0$, the equations
of $s_0$ reduce to $z_0=z_2=\cdots=z_{r-1}=0$ and those of
$s_1$ reduce to $z_1-z_0=z_2=\cdots=z_{r-1}=0$.
Thus the underlying supports are precisely $C\cup D_i$.

At $P_\infty$ use $t=z_r=1$ as frames and affine coordinates.
The tail equations successively eliminate
\[
             z_{r-1}=-s,\quad z_{r-2}=s^2,\quad\ldots,
                         z_2=(-s)^{r-2}.
 \tag{AST22}
\]
For $r=2$ this simply says $z_2=1$ in the chosen chart.
For $Y_0$, its second equation eliminates
$z_0=-s z_2$, and its first remaining equation is $s z_1=0$.
For $Y_1$, its first equation eliminates
$z_0=z_1-s z_2$, and the sum of its first two equations is
again $s z_1=0$.  Hence in each case the local ring after
these eliminations is
\[
                         K[s,z_1]_{(s,z_1)}/(s z_1),
 \tag{AST23}
\]
with the corresponding localizations away from the origin.
This is reduced; its two branches have multiplicity one.
The branch $z_1=0$ is $C$, while the branch $s=0$ is $D_i$.
The eliminated $r-1$ linear equations followed by the
non-zero-divisor $s z_1$ are a regular sequence.

Away from $P_\infty$ on $D_0$, take $z_1=t=1$.
The first equation is $s=0$; the second and the tail then
eliminate $z_0,z_2,\ldots,z_{r-1}$, leaving the one coordinate
$z_r$.  Away from $P_\infty$ on $D_1$, the same chart applies
because $z_1=z_0\ne0$.  The sum of its first two equations
is $s=0$, and the other equations eliminate $z_0-1$ and
$z_2,\ldots,z_{r-1}$, again leaving $z_r$.
These local schemes are smooth and reduced.  Together with
(AST23) and the finite-base chart they cover the zero schemes.
Each is a codimension-$r$ local complete intersection in the
smooth ambient variety; its displayed $r$ generators form a
regular sequence.  The explicit reduced local rings have only
the listed minimal components and no embedded components.
This proves (AST21), including all multiplicities.
\end{proof}

Consequently this calibration has the nonzero original cycle
\[
                  Z=[D_1]-[D_0],\qquad
                  C_U=C^{\rm an}\setminus\{P_\infty\}
                                      \simeq\mathbf A^{1,\rm an}_K.
 \tag{AST24}
\]
The cycle is geometrically homologically trivial, as also follows
by subtracting the cycle classes of two regular zero schemes
of the same $L^{\oplus r}$.  The canceled curve accumulates at
the remaining support at $P_\infty$, so it is not closed in
$X^{\rm an}$.  Nevertheless (AST17) contains all of $C_U$,
and the original-support theorem applies.  The two first-column
pivots of (AST15) are $\tau$ and $1-\tau$; each vanishes at a
point of $C_U$.  Thus neither one of the displayed Gaussian
charts covers the cancellation locus by itself.  This checks
the automatic overlap mechanism without using a closed-cancel
separation assumption.  The example has $n=r+1$ in every rank;
it calibrates the mechanism rather than restricting the theorem
to that dimension.

\paragraph{Limits and exact inputs.}
The new conclusion requires both a literally shared tail and
one affine common-ideal neighborhood (AST3).  It is not a theorem
for arbitrary equal ideals on an arbitrary cover, nor for arbitrary
matrices in $\mathrm{GL}_r$ without a supplied factorization.
It does not remove the hypotheses of the frozen OFC closed-cancel
class or modify frozen CI0.  General different-ideal cancellation,
the full lower-middle kernel problem, arbitrary homologically
trivial cycles, and the extra Beilinson--Bloch height axioms remain
open here.  No high-order form descent or norm-root choice is used.

GB: \path{f4f87d6aabfc325f7d2f01d629a6f4a9fa80227cb9600cf592ea380b0b5104d9}.
ICC, only its elementary conormal relation lemma:
\path{6836e2a945bfbf49faf41bef2d31a9d11c2af1434ada08048616069b906f4e49}.
SHT: \path{a7bc2d335e69cf81800e97d9433fc563108e548a1f1ac3eb15aac15a0022631f}.
RSP: \path{fbd8a7c32e7ec19e2f6775d48f6167f1322e7e13ab03d931a4305381bb3270b2}.
CI: \path{a7a37d03f46437047aaf3cfe9227ae18ba3bf627b65c6dc2ea5d22f610379033}.
All further exact native relation and partition dependencies are
those locked in GB.  Internal exact-hash review is not formal
verification, external peer review, or proof of external novelty.
